\chapter{Percutaneous Nephrolithotomy}

\section*{Overview}
Percutaneous nephrolithotomy (PCNL) is a procedure that removes large or complex kidney stones through a small incision and tract made directly into the kidney.

\section*{Alternative Names}
PCNL; percutaneous stone removal; nephrolithotripsy

\section*{Principle}
A tract is created from the skin into the renal collecting system, a nephroscope visualizes the stone, and ultrasound, pneumatic, or laser energy fragments it so fragments can be extracted.

\section*{History}
PCNL was introduced in the 1970s as an alternative to open surgery for large kidney stones and remains the standard for stones larger than 2 cm.

\section*{Why It Is Done}
Performed for kidney stones larger than 2 cm, staghorn calculi, stones that fail shock wave lithotripsy, and stones in abnormal kidneys.

\section*{Is This a Primary or Secondary Test or Procedure}
Primary treatment for large and complex renal stones.

\section*{How It Is Performed}
Performed under general anesthesia. A needle and guidewire create the tract under fluoroscopy or ultrasound, it is dilated, the nephroscope enters the kidney, and the stone is fragmented and removed. A nephrostomy tube or stent may be placed.

\section*{Preparation}
Preoperative imaging, urine culture and treatment of infection, and anticoagulation management.

\section*{Contraindications and Precautions}
Active untreated urinary infection, bleeding diathesis, and pregnancy are contraindications.

\section*{Risks}
Bleeding, perforation of the collecting system, infection or sepsis, injury to adjacent organs, and residual stone fragments.

\section*{Aftercare and Recovery}
The nephrostomy tube is removed after drainage clears; imaging confirms residual stones.

\section*{Results and Interpretation}
Postoperative imaging and stone analysis guide further therapy and metabolic evaluation.

\section*{Expected Results}
Complete stone clearance with no significant residual fragments.

\section*{Follow-up}
Metabolic stone evaluation and imaging surveillance; additional procedures may be needed for residual stones.

\section*{References}
\begin{enumerate}
  \item MedlinePlus. Percutaneous nephrolithotomy. U.S. National Library of Medicine; 2025.
  \item Current Medical Diagnosis and Treatment. McGraw-Hill; 2025.
\end{enumerate}

\clearpage
